%%%%%%%%%%%%%%%%%%%%%%%%%%%%%%%%%
%
%       EXERCÍCIO
%
%%%%%%%%%%%%%%%%%%%%%%%%%%%%%%%%%

\ifdefstring{\atividade}{prova}{%
    \ifdefstring{\modo}{objetivo}{%
        \renewcommand{\valorquestao}{\ValorQObj\ ponto}
    }{%
        \renewcommand{\valorquestao}{\ValorQDisc\ pontos}
    }%
}%

\begin{exercicioBanco}[\valorquestao]
Assinale a afirmação verdadeira.

% Define as alternativas
\newcommand{\alternativas}{%
    \begin{center}
        \begin{tabularx}{\textwidth}{X}
            (a) \((\sqrt{3}+1)(\sqrt{3}-1)\) é irracional e \(0,999\cdots\) é racional. \\[5pt]
            (b) \((\sqrt{3}+1)(\sqrt{3}-1)\) é racional e \(0,999\cdots\) é racional. \\[5pt]
            (c) \((\sqrt{3}+1)(\sqrt{3}-1)\) é racional e \(0,999\cdots\) é irracional. \\[5pt]
            (d) \((\sqrt{3}+1)(\sqrt{3}-1)\) é irracional e \(0,999\cdots\) é irracional. \\[5pt]
            (e) \((\sqrt{3}+1)(\sqrt{3}-1)\) e \(0,999\cdots\) não são reais.
        \end{tabularx}
    \end{center}
}

% Define a resposta correta
\newcommand{\resposta}{B}

% Lógica condicional para exibição
\ifdefstring{\atividade}{lista}{%
    \alternativas
    \vspace{0.5em}
    
    \noindent\textbf{Resposta:} letra \textbf{\resposta}.
}{%
    \ifdefstring{\modo}{objetivo}{%
        \alternativas
    }{%
        % modo = discursiva → não mostra alternativas
    }
}
\end{exercicioBanco}

